% =============================================================================
% Chapter 9: Fleet Efficiency — ML Systems Lecture Slides (Vol II)
% =============================================================================
% Built around the chapter thesis:
%   "Useful fleet output is capacity times scheduling, runtime, and program goodput."
% One running job (Llama 2 70B on 128 H100s of a shared pool) is charged at
% each boundary: S = 0.65, R = 0.90, P = 0.48, MPG = 0.28, \$7.09 per useful
% GPU-hour. The deck then shows why the program lever is worth the most even
% though scheduling loses the largest share, and how trades are judged.
%
% Slides are written to be read without a speaker: each frame states its claim
% in the title, shows the evidence, and closes with a bottom line. Speaker notes
% add pacing, questions, and misconceptions.
%
% Numbers are the rendered chapter's computed values (feat/vol2-fleet-efficiency
% at a0b79b884b). Figures: images/make_figures.py.
% =============================================================================
\documentclass[aspectratio=169, 12pt]{beamer}
\usepackage{../../assets/beamerthememlsys}

\mlsyssetup{
  volume       = {Volume II},
  chapter      = {Chapter 9},
  logo         = {../../assets/img/logo-mlsysbook.png},
  instlogo     = {../../assets/img/logo-harvard.png},
  chaptertitle = {Fleet Efficiency},
}

% --- Fonts ---
\usepackage[T1]{fontenc}
\usepackage[scaled=0.9]{helvet}
\usepackage{courier}
\renewcommand{\familydefault}{\sfdefault}

% --- Packages ---
\usepackage{booktabs}
\usepackage{amsmath}
\usepackage{array}
\newcolumntype{L}[1]{>{\raggedright\arraybackslash}p{#1}}

% --- Image paths ---
\graphicspath{
  {images/}
}

% --- Chapter-specific macros ---
\newcommand{\Sg}{G_{\text{sched}}}
\newcommand{\Rg}{G_{\text{run}}}
\newcommand{\Pg}{G_{\text{prog}}}
\newcommand{\MPG}{\text{MPG}}
% Bottom line: the one sentence a reader without a speaker should keep.
\newcommand{\bottomline}[1]{%
  \vfill
  \begin{mlsyscard}[top=0.15em,bottom=0.15em,fontupper=\footnotesize]{crimson}\textbf{Bottom line.} #1\end{mlsyscard}%
}
\definecolor{bookblue}{HTML}{006395}
\definecolor{bookorange}{HTML}{E67817}
\definecolor{bookgreen}{HTML}{008F45}
\definecolor{bookviolet}{HTML}{7E317B}
\definecolor{bookbrown}{HTML}{78492A}
\definecolor{bookred}{HTML}{CB202D}
% Compact reference entry for a long reference list
\newcommand{\mlsysrefsmall}[2]{{\scriptsize\textcolor{crimson}{\textbf{[#1]}} #2}\\[0.01cm]}

% --- Helper: safe image include ---
\newcommand{\safeimg}[2][width=\textwidth,keepaspectratio]{%
  \IfFileExists{images/#2}{\includegraphics[#1]{#2}}{%
    \IfFileExists{#2}{\includegraphics[#1]{#2}}{%
      \fbox{\parbox[c][2.5cm][c]{0.85\linewidth}{%
        \centering\footnotesize\textcolor{midgray}{[Missing image]}}}%
    }%
  }%
}

% --- Section count (must match actual \section{} count) ---
\setcounter{mlsystotalsections}{9}

\title{Fleet Efficiency}
\author{Vijay Janapa Reddi}
\institute{Harvard University}
\date{}

\begin{document}

% =============================================================================
% TITLE SLIDE
% =============================================================================
\mlsystitle{Fleet Efficiency}{Useful fleet output is capacity times\\scheduling, runtime, and program goodput.}{cover_performance_engineering_blueprint.png}

% =============================================================================
% LEARNING OBJECTIVES
% =============================================================================
\begin{frame}{Learning Objectives}
\note{
[1 min]
% -- NARRATE
The lecture follows one training job from the capacity its pool offers to the
useful work it produces. Each objective is one boundary or one decision.
% -- FLEX
[CORE] Return to this list at the wrap-up.
}

\footnotesize
By the end of this lecture you should be able to:
\begin{enumerate}\setlength{\itemsep}{1pt}
  \item \textbf{Explain} how useful output decomposes into \textbf{scheduling, runtime, and program goodput}, with each loss charged to exactly one factor.
  \item \textbf{Measure} goodput under a reproducible \textbf{measurement contract} and localize the lost factor from scheduler records, job timelines, and step traces.
  \item \textbf{Attribute} scheduling and runtime losses using the owners' results.
  \item \textbf{Measure} program goodput with \textbf{MFU} (or \textbf{MBU} for decode) and find the Fleet Law term a fix must shorten.
  \item \textbf{Select} fleet-level program levers: compiler partitioning, overlap under contention, and compilation.
  \item \textbf{Evaluate} a trade between factors by its effect on the \textbf{product}, and place it on the \textbf{efficiency frontier}.
\end{enumerate}
\end{frame}

% =============================================================================
% SECTION 1: CAPACITY GAP
% =============================================================================
\section{Capacity Gap}

\begin{frame}{A busy fleet can turn most of its capacity into nothing}
\note{
[3 min]
% -- LINK
Earlier chapters each removed one kind of loss, in their own units:
communication, parallelism, fault tolerance, orchestration. None says how much
capacity remains after all four, or which loss to fix first.
% -- NARRATE
Llama 3 405B is one of the best-documented runs on record. Its two famous
numbers measure different things: more than 90\% effective training time
(reliability) and about 41\% MFU (the speed of the steps it kept). Multiply:
about 37\% of the capacity it held became progress at the ideal rate.
% -- ENGAGE
Ask before revealing the product: ``If 90\% of the time was effective, what
fraction of the capacity became progress?'' Most will say 90\%.
% -- WARN
The two numbers come from different windows, so the product is an estimate.
% -- FLEX
[CORE]
}

\begin{columns}[T]
\begin{column}{0.52\textwidth}
\small
\textbf{Llama 3 405B} on 16,384 H100s:
\begin{itemize}\setlength{\itemsep}{2pt}
  \item more than \textbf{90\%} of elapsed time was effective training time, despite hundreds of interruptions;
  \item the steps it kept ran at about \textbf{41\%} of peak (MFU).
\end{itemize}
\vspace{0.1cm}
$0.90 \times 0.41 \approx$ \textbf{0.37} of held capacity became progress at the ideal rate.
\end{column}
\begin{column}{0.44\textwidth}
\begin{mlsyscard}{bookred}
\footnotesize A dashboard counts every held accelerator as busy. It cannot see a gang still waiting for members, a job replaying lost steps, or a step that runs at half the hardware's rate.
\end{mlsyscard}
\end{column}
\end{columns}

\bottomline{Most of the loss in this well-run job was in slow steps, not in the failures that dominate its reputation.}
\end{frame}

\begin{frame}{Capacity is lost in three places along a job's life}
\note{
[3 min]
% -- NARRATE
Walk left to right. A job must first hold all its devices: until then nothing
runs. Then it must keep its progress: startup, input waits, checkpoint
pauses, failures, and rework consume held time. Then its kept steps must run
at the rate the hardware allows.
% -- ENGAGE
Ask: ``Where does a gang holding half its devices lose capacity?'' [Before
the first gate: it cannot take a single step.]
% -- FLEX
[CORE]
}

\centering
\safeimg[width=0.98\textwidth,height=4.0cm,keepaspectratio]{capacity-funnel.pdf}

\raggedright\footnotesize
\textbf{Running job:} Llama 2 70B on 128 H100s of a shared pool (TP 8 inside each node, DP across 16 nodes). \textbf{Serving twin:} one 8-GPU node per replica.

\bottomline{Scheduling goodput, runtime goodput, and program goodput name the three places, in the order a job meets them.}
\end{frame}

% =============================================================================
% SECTION 2: GOODPUT
% =============================================================================
\section{Goodput}

\mlsysfocus{ML Productivity Goodput}{%
$\MPG = \Sg \times \Rg \times \Pg$\\[0.4cm]
\normalsize scheduling $\times$ runtime $\times$ program: the share of fleet capacity that becomes useful work at the ideal rate%
}

\begin{frame}{Each factor divides by what the previous one passed}
\note{
[4 min]
% -- NARRATE
Read the three ratios down the table. Each numerator is the next denominator,
so they cancel to ideal-rate useful time over offered capacity. That is why
the product is a fraction of capacity, not a rate. The source is Google's
TPU fleet accounting (Wongpanich et al., 2025).
% -- WARN
The classic error: multiplying a run-averaged MFU, which already includes
restart time, by a runtime factor that charges the same restarts again.
% -- FLEX
[CORE]
}

\small
\begin{tabular}{@{}L{0.17\textwidth}L{0.42\textwidth}L{0.33\textwidth}@{}}
  \toprule
  \textbf{Factor} & \textbf{Ratio} & \textbf{Loses capacity to} \\
  \midrule
  \textcolor{bookbrown}{\textbf{Scheduling} $\Sg$} & time held by complete allocations $\div$ fleet capacity-time & partial gangs, fragmentation, idle capacity \\
  \textcolor{bookgreen}{\textbf{Runtime} $\Rg$} & retained-step time $\div$ held time & startup, input waits, checkpoint pauses, failures, rework \\
  \textcolor{bookorange}{\textbf{Program} $\Pg$} & ideal time $\div$ actual time of retained steps & exposed communication, stragglers, stalls, dispatch \\
  \bottomrule
\end{tabular}

\vspace{0.1cm}
\footnotesize Numerators cancel denominators: what remains is \textbf{ideal-rate useful time $\div$ offered capacity}. Useful output $=$ capacity $\times$ MPG; cost per useful unit $=$ price $\div$ MPG.

\bottomline{The product is valid only when no lost interval is subtracted twice or missed between two factors.}
\end{frame}

\begin{frame}{Loss shares add up because each interval is charged once}
\note{
[3 min]
% -- NARRATE
The identity splits the whole gap: the scheduling loss, then the runtime loss
on what was held, then the program loss on what was retained. Apply it to
Llama 3: S = 1 (dedicated cluster), R about 0.90, P about 0.41.
% -- ENGAGE
Ask: ``If every interruption were removed, how much would come back?'' [At
most a tenth; the slow steps hold more than half.]
% -- FLEX
[CORE]
}

\centering
{\normalsize $1 - \Sg\Rg\Pg = (1-\Sg) + \Sg(1-\Rg) + \Sg\Rg(1-\Pg)$}

\vspace{0.15cm}
\safeimg[width=0.98\textwidth,height=3.2cm,keepaspectratio]{llama3-split.pdf}

\raggedright\footnotesize
Runtime boundary $1 - 0.90 = 0.10$; program boundary $0.90 \times 0.59 \approx 0.53$; useful $0.90 \times 0.41 \approx 0.37$.

\bottomline{For one of the most carefully engineered runs on record, the program boundary, not reliability, held most of the gap.}
\end{frame}

\begin{frame}{Program goodput needs a named ideal: MFU, or MBU for decode}
\note{
[4 min]
% -- NARRATE
Activity counters cannot measure program goodput: a device waiting on a
collective looks as busy as one doing math. For training and prefill the
ideal is the required FLOPs at peak, so P is MFU over retained steps. For
batch-one decode the replica must stream its whole weight shard per token:
about 5.2 ms per token, with arithmetic at about 0.3\% of peak. Charging that
to the program would blame it for a loss only batch size can fix, so the
ideal is the bandwidth floor and P is MBU.
% -- WARN
HFU (counting recomputation) exceeds MFU; the gap is recomputation, which is
a program loss even when it is the right price for memory.
% -- FLEX
[CORE]
}

\begin{columns}[T]
\begin{column}{0.5\textwidth}
\small
\textbf{Training and prefill:} the ideal is the model's required FLOPs at peak.
\[
\Pg = \frac{F_{\text{useful}}/(N\,\Rpeak)}{T_{\text{step}}} = \text{MFU}
\]
measured over \textbf{retained steps only}. A run-averaged MFU already includes restarts.
\end{column}
\begin{column}{0.46\textwidth}
\small
\textbf{Batch-one decode:} the ideal is the \textbf{bandwidth floor}, the time to read the required bytes at peak memory bandwidth.

\vspace{0.1cm}
Serving replica: floor $\approx$ \textbf{5.2 ms/token}, arithmetic at about \textbf{0.3\%} of peak. Program goodput is \textbf{MBU}.
\end{column}
\end{columns}

\bottomline{High MBU with low MFU marks a memory-saturated program, not an inefficient one.}
\end{frame}

\begin{frame}{Every loss is charged at the first boundary it fails to cross}
\note{
[3 min]
% -- NARRATE
The boundary decides where a loss is counted; the cause decides who fixes it.
A straggler is a reliability problem, but its cost is a program loss: the job
holds a complete allocation and keeps every step, it just runs them slowly.
A synchronous checkpoint pause is runtime; an asynchronous write that
contends with the step is program.
% -- FLEX
[CORE]
}

\footnotesize
\begin{tabular}{@{}L{0.42\textwidth}L{0.20\textwidth}L{0.31\textwidth}@{}}
  \toprule
  \textbf{Loss} & \textbf{Charged to} & \textbf{Lever usually owned by} \\
  \midrule
  gang holding part of its devices, waiting & scheduling & orchestration (atomic admission) \\
  fragments no queued gang fits & scheduling & orchestration (bin packing) \\
  zombie job holding a complete allocation & runtime & fault tolerance, leases \\
  input wait inside a retained step & runtime & data storage (input pipeline) \\
  preemption: lost work and reload & runtime & orchestration (budgets) \\
  straggling rank slowing every step & program & fault tolerance (replacement) \\
  exposed collectives, noisy neighbor & program & collectives, layout, isolation \\
  \bottomrule
\end{tabular}

\bottomline{Charging a loss to the team that fixes it breaks the product; keep the factor and the owner as two columns.}
\end{frame}

\begin{frame}{Your turn: which factor does each loss lower?}
\note{
[3 min] Exercise. Cover the previous slide.
% -- NARRATE
Two minutes. Then reveal. Item 4 is where most students go wrong: the wait
is in the queue, but the capacity was held by other jobs, so the pool lost
nothing; the job only lost time to result.
% -- FLEX
[CORE]
}

\centering
{\large Name the factor each loss lowers.}\\[0.25cm]
\raggedright\small
\begin{enumerate}\setlength{\itemsep}{3pt}
  \item The job compiles for 12 minutes after its gang completes, before its first step.
  \item An asynchronous checkpoint write slows every step by 3\%.
  \item One rank runs 10\% slow for a week; every synchronous step waits for it.
  \item The job waits 4 hours in the queue while other jobs hold every device.
\end{enumerate}

\pause
\vspace{0.1cm}
\begin{mlsyscard}{crimson}
\small (1) runtime: startup. (2) program: contention inside retained steps. (3) program: every step slowed, though reliability decides replacement. (4) none of the three: the pool lost nothing; the wait costs the job time to result.
\end{mlsyscard}
\end{frame}

\begin{frame}{``Goodput'' names four different measures}
\note{
[2 min]
% -- NARRATE
Every efficiency claim must name the measure, the boundary it divides by,
and the window. SLO goodput applies after MPG: work served at a high average
rate can still miss its latency targets in the tail. Statistical goodput
can fall while program goodput rises when a larger batch makes each example
worth less.
% -- FLEX
[OPTIONAL]
}

\footnotesize
\begin{tabular}{@{}L{0.24\textwidth}L{0.70\textwidth}@{}}
  \toprule
  \textbf{Measure} & \textbf{What it counts} \\
  \midrule
  Run-level goodput & one job's useful training progress after wasted and repeated work \\
  ML productivity goodput & fraction of fleet capacity that becomes useful work at the ideal rate (this lecture) \\
  SLO goodput & rate of completed requests that meet their latency targets (serving) \\
  Statistical goodput & throughput discounted by the optimization progress each example contributes \\
  \bottomrule
\end{tabular}

\bottomline{A goodput number without its measure, boundary, and window cannot be compared with another.}
\end{frame}

% =============================================================================
% SECTION 3: MEASURING
% =============================================================================
\section{Measuring}

\begin{frame}{A goodput claim is valid only under a measurement contract}
\note{
[3 min]
% -- NARRATE
Two measurements of the same job a week apart can differ because they traced
different ranks, spanned a framework upgrade, began before compilation
finished, or averaged over a restart. The contract fixes what must stay
the same. MLPerf is the working model: fixed workload, quality target, rules.
% -- ENGAGE
Ask: ``Which contract field did the MLPerf series on the right violate if you
used it to judge a model change?'' [System state: software version.]
% -- FLEX
[CORE]
}

\begin{columns}[T]
\begin{column}{0.56\textwidth}
\footnotesize
\begin{tabular}{@{}L{0.33\textwidth}L{0.60\textwidth}@{}}
  \toprule
  \textbf{Field} & \textbf{Fixes} \\
  \midrule
  Workload & model revision, input distributions, batch, quality gate \\
  System state & devices, topology, clocks, software and compiler versions, scheduler policy \\
  Timing boundary & setup inside the window, synchronization, cold vs.\ warm \\
  Statistical report & factor, denominator, window, trials, spread \\
  \bottomrule
\end{tabular}
\end{column}
\begin{column}{0.40\textwidth}
\begin{mlsyscard}{bookviolet}
\footnotesize \textbf{Software alone moves results.} The same 8-GPU DGX H100 in four MLPerf Inference rounds (Sep 2023 to Apr 2025): GPT-J 2.7$\times$, Llama 2 70B 1.4$\times$, Stable Diffusion XL 1.2$\times$; ResNet-50 and BERT unchanged.
\end{mlsyscard}
\end{column}
\end{columns}

\bottomline{Two measurements under one contract differ only in the change being evaluated.}
\end{frame}

\begin{frame}{Each factor loses time at its own rate, so each needs its own window}
\note{
[2 min]
% -- NARRATE
Program losses recur in every step: a window of warmed steps is enough.
Runtime losses are rare and costly: the running job meets a hardware failure
about once every 13 days, so a one-day profile usually contains none.
Scheduling follows demand, which fills queues by day and drains them at night.
% -- WARN
A product of the three is meaningful only when all three describe the same
calendar window.
% -- FLEX
[CORE]
}

\small
\begin{tabular}{@{}L{0.18\textwidth}L{0.36\textwidth}L{0.38\textwidth}@{}}
  \toprule
  \textbf{Factor} & \textbf{How its losses occur} & \textbf{Window that sees them} \\
  \midrule
  Program & in every step & warmed steady-state steps \\
  Runtime & rare, expensive events (the running job: a failure about every 13 days) & days to weeks, spanning many interruptions \\
  Scheduling & with the cycle of demand & whole demand cycles (day and night) \\
  \bottomrule
\end{tabular}

\vspace{0.2cm}
Cold-start and warmed trials are reported separately: compilation before the first step is runtime; the speed of the compiled step is program.

\bottomline{A one-day measurement can show runtime goodput near one for a job that loses a tenth of its time over a month.}
\end{frame}

\begin{frame}{Localize the factor before the mechanism: cheapest record first}
\note{
[3 min]
% -- NARRATE
Each record settles one factor and passes the remainder down. The running
job: scheduler records showed the gang completing late; the timeline showed
startup, checkpoint pauses, rework, and a small data stall, and every
retained step longer than planned; only then were traces taken from one
data-parallel group, showing collectives no longer hidden.
% -- WARN
Starting from a trace can declare a job healthy while a third of its held
time goes to restarts no trace window contained.
% -- FLEX
[CORE]
}

\centering
\safeimg[width=0.98\textwidth,height=4.4cm,keepaspectratio]{localization-ladder.pdf}


\bottomline{Each record answers one question the others cannot; reading them out of order gives correct evidence and a wrong conclusion.}
\end{frame}

\begin{frame}{Fleets detect always and trace selectively}
\note{
[3 min]
% -- NARRATE
A 5-second trace of one GPU at 500 MB becomes about 512 GB across 1,024 GPUs,
and full tracing perturbs the step. So an always-on layer records step time,
MFU, per-rank timers, and telemetry; targeted traces capture short windows
from ranks chosen to test a hypothesis. MegaScale's per-rank heat map found
about 0.5\% of machines running substantially slower than the rest.
% -- ENGAGE
Ask: ``Why trace a suspected straggler and a healthy peer together?'' [The
comparison isolates the slow interval.]
% -- FLEX
[OPTIONAL]
}

\begin{columns}[T]
\begin{column}{0.5\textwidth}
\small
\textbf{Always on, every job:} step time, MFU of retained steps, per-rank step timers, device telemetry. Detects that a factor moved.

\vspace{0.15cm}
\textbf{Targeted, on demand:} short traces from chosen ranks; statistical sampling stratified by job, stage, phase, and hour; triggers on outliers.
\end{column}
\begin{column}{0.46\textwidth}
\begin{mlsyscard}{bookorange}
\footnotesize \textbf{Scale of the problem.} 500 MB per GPU for a 5-second trace becomes about \textbf{512 GB} at 1,024 GPUs.\\[3pt]
\textbf{MegaScale:} per-rank timers on every step found about 0.5\% of machines running slow; a few hosts 10\% slower cost about 0.7\% MFU.
\end{mlsyscard}
\end{column}
\end{columns}

\bottomline{A fraction of a percent of machines can set the step time of every rank synchronized with them.}
\end{frame}

\begin{frame}{Regressions are caught in the record that settles each factor}
\note{
[3 min]
% -- NARRATE
Canaries catch most program regressions within minutes. Runtime and
scheduling regressions need baselines kept per software version and per
pool. Cloudflare (2019): one firewall rule's regex backtracking drove the
CPUs serving HTTP traffic to nearly 100\% across the network for 27 minutes;
traffic fell about 82\%. Correct outputs, unacceptable cost, global rollout.
% -- FLEX
[OPTIONAL]
}

\begin{columns}[T]
\begin{column}{0.56\textwidth}
\footnotesize
\begin{tabular}{@{}L{0.36\textwidth}L{0.17\textwidth}L{0.38\textwidth}@{}}
  \toprule
  \textbf{Tier} & \textbf{Factor} & \textbf{Catches} \\
  \midrule
  small-scale canaries & program & scaling-curve bends \\
  fleet MFU baselines & program & compiler or config changes \\
  per-rank step timers & program & gray failures, stragglers \\
  startup, recovery baselines & runtime & slower compiles, checkpoints \\
  allocation baselines & scheduling & policies that strand capacity \\
  \bottomrule
\end{tabular}
\end{column}
\begin{column}{0.40\textwidth}
\begin{mlsyscard}{bookred}
\footnotesize\textbf{Cloudflare, July 2019.} A firewall rule's regex backtracking drove edge CPUs near 100\%: a 27-minute outage, traffic down about 82\%.
\end{mlsyscard}
\end{column}
\end{columns}

\bottomline{Stage performance-sensitive changes on a small share of the fleet and promote them only when the factors they touch hold.}
\end{frame}

% =============================================================================
% SECTION 4: SCHEDULING
% =============================================================================
\section{Scheduling}

\begin{frame}{Only complete allocations count toward scheduling goodput}
\note{
[3 min]
% -- NARRATE
A rigid job holding half its gang completes zero steps, not half as many, yet
a utilization counter credits every held device. Scheduling goodput divides
complete-allocation time by in-service capacity. Devices out for repair are
reported beside it as availability. Low S describes the pool: free capacity
with an empty queue is a planning question; free capacity with waiting work
is a policy question.
% -- FLEX
[CORE]
}

\small
\[
\Sg = \frac{\sum_{\text{jobs}} \text{devices} \times \text{(complete} \to \text{release)}}{\text{in-service devices} \times \text{window}}
\]

\begin{columns}[T]
\begin{column}{0.48\textwidth}
\footnotesize
\textbf{Counted as loss:} capacity no job holds; gangs holding some members; fragments no queued gang fits; reservations draining for the head of the queue.
\end{column}
\begin{column}{0.48\textwidth}
\footnotesize
\textbf{Not counted:} a job's own queue wait while other jobs hold the devices (it costs time to result, not pool capacity); devices out for repair (reported as availability).
\end{column}
\end{columns}

\bottomline{Queue depth recorded beside free capacity tells a capacity-planning problem from a placement-policy problem.}
\end{frame}

\begin{frame}{A shared pool can lose a third of its capacity before any step runs}
\note{
[3 min]
% -- NARRATE
Constructed snapshot: 80\% assigned, of which 15\% are gangs still waiting.
Scheduling goodput is 0.80 minus 0.15: 0.65. The 5\% of input waits sits
inside complete allocations, so it is a runtime loss. Atomic group admission
returns the 15\% to the pool, but S rises to at most 0.80, and only if other
jobs' complete allocations take every freed device.
% -- ENGAGE
Ask: ``Google reports S above 95\% for every job size. Why are the largest
and smallest jobs highest?'' [Large jobs are protected from preemption;
small ones refit easily and are used to defragment.]
% -- FLEX
[CORE]
}

\centering
\safeimg[width=0.98\textwidth,height=3.9cm,keepaspectratio]{pool-snapshot.pdf}

\raggedright\footnotesize
Loss $= 0.20$ unheld $+ 0.15$ in incomplete gangs $= 0.35$. Group admission alone lifts $\Sg$ to at most 0.80.

\bottomline{Every gain inside a job reaches this pool at 0.65 of its value on a dedicated cluster.}
\end{frame}

% =============================================================================
% SECTION 5: RUNTIME
% =============================================================================
\section{Runtime}

\begin{frame}{Runtime goodput counts held time that becomes retained steps}
\note{
[3 min]
% -- NARRATE
Numerator: retained steps times the job's own fault-free step time.
Denominator: held time from gang completion to release. Normalizing to the
job's own rate keeps program losses out: a straggler that slows every step
for days is a program loss; the same slowdown confined to the minutes after a
restart is a runtime loss.
% -- WARN
A preempted job's wait in the queue before readmission is not held time.
% -- FLEX
[CORE]
}

\small
\[
\Rg = \frac{\text{retained steps} \times \text{fault-free step time}}{\text{held time}}
\]

\textbf{Runtime losses, in the order a job meets them:} startup and compilation before the first step; input waits inside retained steps; synchronous checkpoint pauses; failure detection, relaunch, reload, warmup; rework of steps lost since the last checkpoint; preemption.

\vspace{0.15cm}
\footnotesize\textbf{Charged elsewhere:} queue wait after an eviction (not held); asynchronous checkpoint contention and compiled-code speed (program).

\bottomline{The job's own fault-free rate is the reference, so slow steps stay a program loss and only interruptions are runtime.}
\end{frame}

\begin{frame}{Five owned runtime terms add up}
\note{
[4 min]
% -- NARRATE
Each owning chapter gives its loss as a share of held time. For the running
job over a week: input waits 7.7\% (the snapshot's 5\% of capacity over S =
0.65), preemption 1.4\% (1 per day at 20.2 min), checkpoint 0.4\% (9.8 s
writes at a 77-min interval plus rework), startup 0.05\%, recovery 0.03\%
(MTBF about 302 h). Sum about 10\%: R about 0.90.
% -- WARN
Each interval enters one term only: rework sits in the checkpoint term, not
again in recovery; warmup belongs to the event that caused it.
% -- FLEX
[CORE]
}

\begin{columns}[T]
\begin{column}{0.52\textwidth}
\safeimg[width=\textwidth,height=4.3cm,keepaspectratio]{runtime-terms.pdf}
\end{column}
\begin{column}{0.45\textwidth}
\footnotesize
\[
1 - \Rg \approx \frac{t_{\text{start}}}{t_{\text{held}}} + f_{\text{stall}} + f_{\text{ckpt}} + \frac{t_{\text{recover}}}{\text{MTBF}} + \lambda_{\text{pre}}\,t_{\text{pre}}
\]
Llama 2 70B, 128 GPUs, one week, one forced preemption a day. Input waits and preemption are 9.1 of the 10 points; checkpoint and recovery together are under half a point.
\end{column}
\end{columns}

\bottomline{On a mid-sized shared job, the input pipeline reaches about 17$\times$ more held time than faster checkpoints.}
\end{frame}

\begin{frame}{Job size and length reorder the runtime terms}
\note{
[3 min]
% -- NARRATE
Checkpoint waste and recovery grow with device count; input waits, startup,
and preemption do not. At the 10,000-GPU reference, checkpoint waste (about
5.9\%) and recovery (about 2.8\%) take nearly a tenth of held time; the two
pass input waits near 13,000 GPUs. Startup is negligible for a week-long job
and about 9\% of an hour-long evaluation job.
% -- ENGAGE
Ask: ``Counting the preemption budget as interruptions, what happens to the
Young-Daly interval?'' [77 min shortens to about 21 min.]
% -- FLEX
[OPTIONAL]
}

\small
\begin{tabular}{@{}L{0.42\textwidth}L{0.50\textwidth}@{}}
  \toprule
  \textbf{Situation} & \textbf{Dominant runtime term} \\
  \midrule
  running job, 128 GPUs, one week & input waits 7.7\%, preemption 1.4\% \\
  10,000-GPU reference & checkpoint 5.9\% + recovery 2.8\% \\
  12 forced preemptions a day & preemption 16.8\% \\
  one-hour evaluation job & startup about 9\% \\
  \bottomrule
\end{tabular}

\vspace{0.2cm}
\textbf{Coupling no single owner sees:} a preemption discards work exactly as a failure does, so the preemption budget changes the checkpoint interval (77 min to about 21 min).

\bottomline{The same runtime goodput can point to different owners, so report the terms, not only their sum.}
\end{frame}

\begin{frame}{Discussion: where would you invest to raise runtime goodput?}
\note{
[4 min] Discussion in pairs.
% -- NARRATE
Expected: for the 128-GPU job, the input pipeline and the preemption budget,
not checkpoint speed. For a 16,000-GPU dedicated run, checkpoint and recovery.
For a sweep of hour-long jobs, a shared compilation cache to cut startup.
% -- WARN
Students will default to faster checkpointing because it is the famous lever.
% -- FLEX
[CORE]
}

\centering
\vspace{0.2cm}
{\large Your reliability team proposes faster checkpoint writes and warm restart.}\\[0.3cm]
\raggedright\small
\begin{enumerate}\setlength{\itemsep}{3pt}
  \item For the running job (128 GPUs, one week), what is the most this can recover?
  \item Which lever reaches the most held time instead, and by how much more?
  \item For which kind of job would the reliability team's proposal be the right first move?
\end{enumerate}
\centering
{\small\textcolor{midgray}{Two minutes with a neighbor.}}
\end{frame}

% =============================================================================
% SECTION 6: PROGRAM
% =============================================================================
\section{Program}

\begin{frame}{A local fix reaches the fleet only through the term it shortens}
\note{
[3 min]
% -- LINK
The single-machine iron law splits a device's time into data, compute, and
overhead. The Fleet Law adds communication and synchronization across
devices and credits back what overlaps with useful work.
% -- NARRATE
A kernel speedup shortens local execution only. Local execution divides
across devices while communication and synchronization do not shrink in
proportion, so the same kernel win is worth less to a larger job.
% -- FLEX
[CORE]
}

\footnotesize
\begin{tabular}{@{}L{0.21\textwidth}L{0.30\textwidth}L{0.41\textwidth}@{}}
  \toprule
  \textbf{Program loss} & \textbf{Term it lengthens} & \textbf{Evidence in a step trace} \\
  \midrule
  memory stalls & $D/BW$ (local) & kernels near the bandwidth ceiling \\
  arithmetic below peak & $O/(\Rpeak\,\eta)$ (local) & kernels far below both ceilings \\
  recomputation & added $O$ & HFU above MFU \\
  host dispatch & dispatch part of $L$ & gaps between short kernels \\
  exposed communication & $T_{\text{comm}} - T_{\text{overlap}}$ & collectives with no concurrent compute \\
  stragglers, pipeline bubbles & $T_{\text{sync}}$ & all ranks but one waiting at a barrier \\
  \bottomrule
\end{tabular}

\bottomline{Name the term before choosing a lever; a fix for a term that is not exposed buys nothing.}
\end{frame}

\begin{frame}{Scale cost the running job a quarter of its node efficiency}
\note{
[3 min]
% -- NARRATE
Same work per device in both configurations. Node: 180 ms per step, MFU 65\%.
128 GPUs: 245 ms, MFU 48\% (14.9 PFLOP useful over the 31 PFLOP the GPUs
could execute in 245 ms). Ratio 180/245 is about 74\%: the scaling efficiency
from one node.
% -- WARN
Report fleet MFU, or node MFU and scaling efficiency, never fleet MFU times
scaling efficiency: that charges the exposed time twice.
% -- FLEX
[CORE]
}

\centering
\safeimg[width=0.95\textwidth,height=3.5cm,keepaspectratio]{node-vs-fleet.pdf}

\raggedright\footnotesize
Fleet MFU $= 65\% \times 0.74 \approx 48\%$ $=$ program goodput. The 65 ms beyond the node bounds what exposed communication and synchronization can claim; a trace decides how it divides. Llama 3 405B ran at about 41\% on 16,384 GPUs.

\bottomline{The 65 ms per step beyond one node is the program loss that scale adds; only a step trace attributes it.}
\end{frame}

\begin{frame}{Drill down until a level names the term}
\note{
[3 min]
% -- NARRATE
Each level is taken only for the ranks and window the level above selected.
The running job stopped at kernel timelines: no rank consistently behind
(not a straggler), collectives on their own stream with no computation
beside them during part of the backward pass. Counters were not needed.
% -- ENGAGE
Ask: ``LayerNorm achieves 5 TFLOP/s and 2.8 TB/s on an H100. Which lever?''
[1.8 FLOP/byte, far below the 295 ridge, at 84\% of peak bandwidth: move
fewer bytes by fusing, not faster math.]
% -- FLEX
[CORE]
}

\begin{columns}[T]
\begin{column}{0.48\textwidth}
\safeimg[width=\textwidth,height=4.3cm,keepaspectratio]{drill-hierarchy.pdf}
\end{column}
\begin{column}{0.49\textwidth}
\small
\textbf{Running job:} no lagging rank; collectives on their own stream with nothing beside them. The loss is exposed communication, a schedule problem, not bandwidth.

\vspace{0.15cm}
\footnotesize\textbf{Roofline reading:} a LayerNorm at 5 TFLOP/s and 2.8 TB/s has intensity 1.8 FLOP/byte, far below the H100 ridge of about 295, at 84\% of peak bandwidth. Only moving fewer bytes helps.
\end{column}
\end{columns}

\bottomline{A busy device is placed on the roofline; an idle one is a wait on input, dispatch, or communication first.}
\end{frame}

\begin{frame}{Three levers reach three different terms}
\note{
[4 min]
% -- NARRATE
Compiler-derived partitioning (GSPMD): annotate a few tensors along named
mesh axes; the compiler propagates layouts and inserts the collectives. An
unexpected AllGather in every layer is the signature of a layout conflict.
Overlap hides part of the communication, but streams share SMs and memory
bandwidth: it pays only if the time removed from the critical path exceeds
the time added to compute. Kernel levers act on local execution only.
% -- FLEX
[CORE]
}

\scriptsize
\begin{tabular}{@{}L{0.22\textwidth}L{0.27\textwidth}L{0.43\textwidth}@{}}
  \toprule
  \textbf{Lever} & \textbf{Term it moves} & \textbf{Evidence it worked} \\
  \midrule
  compiler-derived partitioning & chooses $T_{\text{comm}}$ and the chance to overlap & collectives in the trace match what the layout predicts \\
  overlap under contention & raises $T_{\text{overlap}}$ & collective intervals move under compute, and the step gets shorter \\
  fusion (about 1.7$\times$ less HBM traffic) & $D/BW$, dispatch & fewer kernels, fewer HBM bytes \\
  tiled attention (about 65$\times$ at 8,192 tokens) & $D/BW$, HBM capacity & no score-sized allocations \\
  graph replay & dispatch only & launch gaps vanish (eager dispatch is 12--24 ms per 70B forward) \\
  \bottomrule
\end{tabular}

\bottomline{The running job's largest program loss, the data-axis gradient exchange, is out of reach of every kernel lever.}
\end{frame}

\begin{frame}{Model-axis exchanges cost the node; data-axis, the fleet}
\note{
[3 min]
% -- NARRATE
Model-axis AllReduces complete activations the next operation needs: about
36 MB each, 320 per step, about 45 ms (25\% of the 180 ms node step) by a
bandwidth-only estimate. They never leave the node, so they are part of the
node baseline. The data-axis gradient exchange can be bucketed and overlapped
with backward; most of the 65 ms the fleet adds is its exposed part. The
serving replica's per-layer exchange is 32.8 KB: latency-bound, with an
assumed launch and protocol cost about 70$\times$ the transfer.
% -- FLEX
[OPTIONAL]
}

\begin{columns}[T]
\begin{column}{0.48\textwidth}
\begin{mlsyscard}{bookorange}
\footnotesize\textbf{Model axis (inside the node)}\\[2pt]
about 36 MB per exchange, 320 per step, about 45 ms: 25\% of the node step. Strict dependency; does not grow with the fleet.
\end{mlsyscard}
\end{column}
\begin{column}{0.48\textwidth}
\begin{mlsyscard}{bookviolet}
\footnotesize\textbf{Data axis (across nodes)}\\[2pt]
gradient buckets can reduce while backward continues. Its exposed part is most of the 65 ms the fleet step adds.
\end{mlsyscard}
\end{column}
\end{columns}

\vspace{0.15cm}
\footnotesize\textbf{Serving replica:} each per-layer exchange is about 32.8 KB, far below NVLink's bandwidth regime, so launch cost sets its time. Fewer exchanges or graph replay help; overlap cannot.

\bottomline{The first experiment for the running job is overlapping the data-axis exchange, not tuning kernels.}
\end{frame}

\begin{frame}{Program speedups do not add, so re-measure after every fix}
\note{
[2 min]
% -- NARRATE
Reinforcing: fusion and tiling free memory, allowing a larger microbatch,
which lengthens the compute that hides each collective. Redundant: a
compiler that already fuses a chain leaves nothing for a hand fusion.
Conflicting: a faster backward pass shrinks the window that hid the gradient
exchange, exposing communication that was free before.
% -- FLEX
[CORE]
}

\small
\begin{columns}[T]
\begin{column}{0.31\textwidth}
\begin{mlsyscard}{bookgreen}
\footnotesize\textbf{Reinforcing}\\ freed memory allows a larger microbatch, which hides more communication
\end{mlsyscard}
\end{column}
\begin{column}{0.31\textwidth}
\begin{mlsyscard}{midgray}
\footnotesize\textbf{Redundant}\\ the compiler already fused the chain a hand-written kernel targets
\end{mlsyscard}
\end{column}
\begin{column}{0.31\textwidth}
\begin{mlsyscard}{bookred}
\footnotesize\textbf{Conflicting}\\ a faster backward pass shrinks the window that hid the gradient exchange
\end{mlsyscard}
\end{column}
\end{columns}

\vspace{0.25cm}
\textbf{Order:} measure under the contract; confirm the loss is in the program; name the term; apply the cheapest lever for it; re-measure; repeat.

\bottomline{Only an end-to-end measurement classifies a combination of fixes as reinforcing, redundant, or conflicting.}
\end{frame}

% =============================================================================
% SECTION 7: WHOLE FLEET
% =============================================================================
\section{Whole Fleet}

\begin{frame}{The running job: 43\% of held hours, 28\% of the pool's capacity}
\note{
[3 min]
% -- NARRATE
Job product R times P: 0.90 times 0.48, about 0.43, which the job's team can
change. Pool product S times R times P: about 0.28, which the fleet owner
pays for. Loss shares by the identity: 35\% scheduling, 6\% runtime, 31\%
program, 28\% useful. Llama 3 405B moved onto this pool: 0.65 times 0.37,
about 0.24.
% -- WARN
The pool figure assumes every job runs like this one; idle reservations and
zombies make the real figure lower.
% -- FLEX
[CORE]
}

\centering
\safeimg[width=0.95\textwidth,height=2.9cm,keepaspectratio]{pool-loss-shares.pdf}

\raggedright\small
\textbf{Job product} $\Rg\Pg \approx 0.43$: what the job's team can change.\quad
\textbf{Pool product} $\MPG \approx 0.28$: what the fleet owner pays for.

\bottomline{No engineering inside a job recovers a loss that occurs before the job holds its devices.}
\end{frame}

\begin{frame}{Predict: which lever is worth the most?}
\note{
[4 min] Predict, then reveal.
% -- NARRATE
Most students pick scheduling because it lost the largest share. Reveal: a
point of any factor is worth the product of the other two (S: 0.43, R: 0.31,
P: 0.59). At each owner's ceiling, scheduling to 0.80 gives MPG 0.35 (+23\%),
removing input waits (R to 0.98) gives 0.31 (+9\%), program to the node's
0.65 gives 0.38 (+35\%).
% -- WARN
Capacity recovered at the first boundary must still cross the later ones at
their current rates.
% -- FLEX
[CORE]
}

\begin{columns}[T]
\begin{column}{0.48\textwidth}
\small
Scheduling lost \textbf{35\%} of the pool, the program \textbf{31\%}, runtime \textbf{6\%}.

\vspace{0.15cm}
Each owner can reach a ceiling: scheduling to 0.80, runtime to 0.98 (no input waits), program to 0.65 (the node's MFU).

\vspace{0.15cm}
\textbf{Which lever raises MPG the most?}

\pause
\vspace{0.15cm}
\begin{mlsyscard}{crimson}
\footnotesize The program: +35\%. A point of a factor is worth the product of the other two (S 0.43, R 0.31, P 0.59).
\end{mlsyscard}
\end{column}
\begin{column}{0.48\textwidth}
\safeimg[width=\textwidth,height=4.8cm,keepaspectratio]{lever-ceilings.pdf}
\end{column}
\end{columns}
\end{frame}

\begin{frame}{Every factor is a cost multiplier on useful work}
\note{
[2 min]
% -- NARRATE
At the internal fleet rate of \$2 per GPU-hour, the pool pays about \$7.09
for each GPU-hour of ideal-rate progress: 3.5$\times$ the capacity price. A
chargeback by held hours bills the job 2.3$\times$ and leaves the scheduling
loss unbilled. Raising program goodput to its ceiling cuts the cost of useful
work by 26\% without buying an accelerator.
% -- FLEX
[CORE]
}

\centering
{\Large $\text{cost per useful GPU-hour} = \dfrac{\text{price}}{\Sg\,\Rg\,\Pg}$}

\vspace{0.3cm}
\raggedright\small
\begin{tabular}{@{}L{0.55\textwidth}r@{}}
  \toprule
  capacity price (internal fleet rate) & \$2 per GPU-hour \\
  cost per useful GPU-hour, running job's pool & \textbf{\$7.09} (3.5$\times$) \\
  chargeback by held hours charges the job & 2.3$\times$ (scheduling loss unbilled) \\
  program goodput raised to its ceiling & cost per useful unit $-26\%$ \\
  \bottomrule
\end{tabular}

\bottomline{A loss at any boundary raises the price of every useful unit, a training token or a served one, by the same proportion.}
\end{frame}

% =============================================================================
% SECTION 8: TRADES
% =============================================================================
\section{Trades}

\begin{frame}{Judge a trade by the product of the factors' ratios}
\note{
[3 min]
% -- NARRATE
A change that touches several factors multiplies MPG by the product of each
factor's after/before ratio. Coupling: a compiler that shortens the on-device
step raises program goodput and lowers runtime goodput, because startup and
checkpoint time do not shrink. A dashboard showing one factor falling after
an improvement has not shown a regression until the product is read.
% -- WARN
Asymmetry: P from 0.48 to 0.65 is +35\% and would pay even at a 26\% runtime
cost; R at 0.90 has at most 11\% left to gain.
% -- FLEX
[CORE]
}

\footnotesize
\begin{tabular}{@{}L{0.26\textwidth}L{0.12\textwidth}L{0.12\textwidth}L{0.19\textwidth}L{0.19\textwidth}@{}}
  \toprule
  \textbf{Lever} & \textbf{Program} & \textbf{Runtime} & \textbf{Scheduling} & \textbf{MPG} \\
  \midrule
  compiler shortens on-device step & rises & falls & falls if device-bound & rises if device-bound \\
  runtime removes idle or lost time & -- & rises & falls & rises \\
  scheduler shortens partial holding & -- & -- & rises & rises \\
  \bottomrule
\end{tabular}

\vspace{0.2cm}
\small A factor near one has little relative gain left to sell: spending runtime goodput to buy program goodput deserves a measurement; the reverse deserves suspicion.

\bottomline{A change is worth making only when the product of its factor ratios exceeds one.}
\end{frame}

\begin{frame}{Your turn: does this change pay?}
\note{
[3 min] Exercise.
% -- NARRATE
Ratios multiply: 1.1 times 11/12 is about 1.008, so the product rises by
about 0.8\%. Scheduling goodput is unchanged and cancels; you do not need its
value. Follow-up: if the change also required a 6-hour placement wait that
held capacity free, the pool's product would have to include that loss.
% -- FLEX
[CORE]
}

\centering
\vspace{0.2cm}
{\large A change raises program goodput by \textbf{one tenth} of its value\\ and lowers runtime goodput by \textbf{one twelfth} of its value.}\\[0.3cm]
{\normalsize Does MPG rise? Do you need to know scheduling goodput?}

\pause
\vspace{0.4cm}
\raggedright
\begin{mlsyscard}{crimson}
\small $1.10 \times \tfrac{11}{12} \approx 1.008$: MPG rises by about 0.8\%. Scheduling goodput is unchanged, so its ratio is 1 and its value is not needed.
\end{mlsyscard}
\end{frame}

\begin{frame}{A compilation build is repaid over executions, not over time}
\note{
[3 min]
% -- NARRATE
Constructed: a 13B model decodes at 120 tokens/s on one H100 (8.33 ms per
token). Autotuned compilation cuts element-wise time 70\% and launch
overhead 80\%: 0.635 of the step, about 5.3 ms, saving about 3.0 ms per
token. A 15-minute build is repaid after about 300,000 tokens, about 26
minutes of one stream. Shared by 128 replicas: about 12 seconds. Built per
replica: about 32 GPU-hours per rollout, a quarter of a burst-hour replica.
% -- WARN
Recompiles triggered by varying shapes in steady operation are a program
loss (slow steps), not startup.
% -- FLEX
[CORE]
}

\begin{columns}[T]
\begin{column}{0.5\textwidth}
\safeimg[width=\textwidth,height=4.7cm,keepaspectratio]{compile-amortization.pdf}
\end{column}
\begin{column}{0.47\textwidth}
\footnotesize
\textbf{Constructed decode budget}, 13B model, one H100: 8.33 ms per token uncompiled; compiled about 5.3 ms (1.58$\times$).

\vspace{0.1cm}
\textbf{Build:} about 15 min autotuned, charged to runtime (startup).

\vspace{0.1cm}
\textbf{Per-replica builds:} 128 $\times$ 15 min $\approx$ 32 GPU-hours per rollout, about 25\% of a replica added for one burst hour.
\end{column}
\end{columns}

\bottomline{Share compiled artifacts through a persistent cache keyed by model, shapes, mesh, accelerator, and compiler version.}
\end{frame}

\begin{frame}{The product decides for the fleet; time to result decides for the job}
\note{
[3 min]
% -- NARRATE
Scheduling goodput is absent from time to result: the pool's idle capacity
does not slow a job that holds its allocation. More devices shorten held
time only while N times program goodput keeps rising. Elastic example: a job
targeting 512 GPUs can start on 128, wait 4 hours for 512, or start small and
grow to 256. The step-up strategy progresses at about 1.7$\times$ the
immediate run and trails the full allocation by about 38\%.
% -- ENGAGE
Ask: ``Which strategy does a team with a deadline pick, and which does a
fleet with a deep queue prefer?''
% -- FLEX
[CORE]
}

\centering
{\large $T_{\text{result}} = T_{\text{queue}} + \dfrac{T_{\text{ideal}}}{N \times \Rg \times \Pg}$}

\vspace{0.1cm}
\raggedright\footnotesize
\begin{tabular}{@{}L{0.20\textwidth}L{0.72\textwidth}@{}}
  \toprule
  \textbf{Dimension} & \textbf{Typical tension} \\
  \midrule
  goodput & a lever that raises one factor often lowers another \\
  cost per useful unit & cheaper capacity (spot, older accelerators) often carries lower runtime or program goodput \\
  time to result & more devices shorten held time but lower program goodput and can lengthen the queue \\
  \bottomrule
\end{tabular}

\bottomline{Different owners choose different points on one efficiency frontier; a change states which point it aims at.}
\end{frame}

% =============================================================================
% SECTION 9: WRAP-UP
% =============================================================================
\section{Wrap-Up}

\begin{frame}{Fallacies}
\note{
[3 min]
% -- NARRATE
Each fallacy reads one factor, one window, or one kernel as the whole.
% -- FLEX
[CORE]
}

\small
\textbf{Fallacy:} \textit{A job with high MFU is using the fleet well.}\\
The running job's steps run at 0.48, above Llama 3's 0.41, yet scheduling at 0.65 and runtime at 0.90 leave an MPG of 0.28.

\vspace{0.15cm}
\textbf{Fallacy:} \textit{A fleet with high capacity, occupancy, and duty cycle is productive.}\\
Each counts accelerators available, held, or busy; none counts progress.

\vspace{0.15cm}
\textbf{Fallacy:} \textit{More FLOP/s means a proportionally faster step.}\\
Batch-one decode uses about 0.3\% of peak: twice the peak leaves its 5.2 ms floor unchanged.

\vspace{0.15cm}
\textbf{Fallacy:} \textit{Fixing the largest loss first recovers the most useful work.}\\
Scheduling loses 35\%, yet at the owners' ceilings the program lever raises MPG to 0.38 and the scheduling lever to 0.35.
\end{frame}

\begin{frame}{Pitfalls}
\note{
[3 min]
% -- NARRATE
Each pitfall breaks either the accounting or the measurement.
% -- FLEX
[CORE]
}

\small
\textbf{Pitfall:} \textit{Charging each loss to the team that will fix it.}\\
A zombie charged to scheduling subtracts held hours twice; a straggler charged to reliability vanishes from program goodput.

\vspace{0.15cm}
\textbf{Pitfall:} \textit{Judging a job by its single profiling run.}\\
The running job meets a failure about once every 13 days; a one-day profile contains none.

\vspace{0.15cm}
\textbf{Pitfall:} \textit{Judging a job by mean step time.}\\
Intermittent stragglers delay their whole group each episode; per-rank timers expose them.

\vspace{0.15cm}
\textbf{Pitfall:} \textit{Counting retained steps as useful work without checking quality.}\\
Lower precision or larger batches can raise program goodput while needing more steps to reach the same quality.
\end{frame}

\begin{frame}{What were the key ideas?}
\note{
[3 min] Retrieval practice. Collect three answers, compare with the takeaways.
% -- FLEX
[CORE]
}

\centering
\vspace{0.8cm}
{\Large\bfseries Close your notes.}\\[0.5cm]
{\normalsize Write the \textbf{three goodput factors}, the \textbf{record that settles each},\\ and \textbf{why the program lever beat the scheduling lever} for the running job.}\\[0.3cm]
{\normalsize 90 seconds, no peeking.}
\end{frame}

\begin{frame}{Key takeaways}
\note{
[2 min]
% -- FLEX
[CORE]
}

\footnotesize
\begin{itemize}\setlength{\itemsep}{2pt}
  \item \textbf{Goodput:} $\MPG = \Sg \Rg \Pg$; charge each loss once, at the first boundary it fails to cross.
  \item \textbf{Measuring:} fix a contract and each factor's window; read scheduler records, then timelines, then traces.
  \item \textbf{Scheduling:} only complete allocations count; a shared pool can lose a third before any step runs.
  \item \textbf{Runtime:} five owned terms add; on a mid-sized shared job, input waits and preemption dominate.
  \item \textbf{Program:} MFU over retained steps (MBU for decode); a fix helps only through the exposed term it shortens.
  \item \textbf{Whole fleet:} a point of any factor is worth the product of the other two.
  \item \textbf{Trades:} a change pays only if the product of its factor ratios exceeds one.
\end{itemize}
\end{frame}

\begin{frame}{The accounting on one page}
\note{
[1 min] Reference slide.
% -- FLEX
[OPTIONAL]
}

\footnotesize
\begin{tabular}{@{}L{0.26\textwidth}L{0.70\textwidth}@{}}
  \toprule
  \textbf{Quantity} & \textbf{Rule} \\
  \midrule
  ML productivity goodput & $\MPG = \Sg\,\Rg\,\Pg$ \\
  Loss shares & $1 - \MPG = (1-\Sg) + \Sg(1-\Rg) + \Sg\Rg(1-\Pg)$ \\
  Program goodput & MFU over retained steps; MBU against the bandwidth floor for decode \\
  Runtime loss & startup share $+$ data stall $+$ checkpoint waste $+$ recovery$/$MTBF $+$ preemption rate $\times$ cost \\
  Value of one point & the product of the other two factors \\
  Cost per useful unit & price $\div$ MPG \\
  Trade rule & product of factor ratios $> 1$ \\
  Time to result & $T_{\text{queue}} + T_{\text{ideal}}/(N\,\Rg\,\Pg)$ \\
  \bottomrule
\end{tabular}
\end{frame}

\begin{frame}{References}
\note{
[0 min] Reference list for further reading.
% -- FLEX
[OPTIONAL]
}

\mlsysrefsmall{Emer84}{Emer and Clark, ``A Characterization of Processor Performance in the VAX-11/780,'' ISCA 1984.}
\mlsysrefsmall{Williams09}{Williams, Waterman, and Patterson, ``Roofline: An Insightful Visual Performance Model,'' CACM 2009.}
\mlsysrefsmall{Mattson20}{Mattson et al., ``MLPerf Training Benchmark,'' MLSys 2020.}
\mlsysrefsmall{Xu21}{Xu et al., ``GSPMD: General and Scalable Parallelization for ML Computation Graphs,'' 2021.}
\mlsysrefsmall{Chowdhery22}{Chowdhery et al., ``PaLM: Scaling Language Modeling with Pathways,'' 2022.}
\mlsysrefsmall{Wang22}{Wang et al., ``Overlap Communication with Dependent Computation via Decomposition,'' ASPLOS 2023.}
\mlsysrefsmall{Dubey24}{Dubey et al., ``The Llama 3 Herd of Models,'' 2024.}
\mlsysrefsmall{Jiang24}{Jiang et al., ``MegaScale: Scaling LLM Training to More Than 10,000 GPUs,'' NSDI 2024.}
\mlsysrefsmall{Kokolis25}{Kokolis et al., ``Revisiting Reliability in Large-Scale ML Research Clusters,'' HPCA 2025.}
\mlsysrefsmall{Wongpanich25}{Wongpanich et al., ``Machine Learning Fleet Efficiency: Analyzing and Optimizing Large-Scale Google TPU Systems with ML Productivity Goodput,'' 2025.}
\end{frame}

\begin{frame}{Next lecture: from useful work to served requests}
\note{
[1 min]
% -- LINK
Everything today averaged over hours to weeks and assumed work was waiting
whenever a device was ready.
% -- FLEX
[CORE]
}

\centering
\vspace{0.4cm}
{\large Requests arrive unevenly, and each must finish within a latency target.}\\[0.4cm]
{\normalsize Queueing, batching, and cache admission decide how much of a replica's measured capacity becomes completed requests.}\\[0.6cm]
\begin{mlsyscard}{crimson}
\centering\normalsize \textbf{Serving at Scale} starts from this replica's bandwidth floor and its time to become ready, and measures SLO goodput, which applies after MPG.
\end{mlsyscard}
\end{frame}

\end{document}
